% ============================================================
\chapter{Tests, validation et finalisation}
\label{chap:sprint6}
% ============================================================

\section{Introduction}
\label{sec:intro_sprint6}

Ce chapitre présente le \textbf{sixième et dernier sprint} de la
plateforme FinBank. Il est consacré à la vérification du système, à la
correction des anomalies et à la préparation de la livraison finale.

La validation porte sur les services backend, les contrôleurs REST, les
fonctionnalités de sécurité, les processus Flowable, l'assistant
bancaire ainsi que les principales interfaces React. Les parcours
utilisateur sont également vérifiés afin de s'assurer que les
fonctionnalités développées au cours des sprints précédents restent
cohérentes après leur intégration.

\section{Backlog du Sprint}
\label{sec:backlog_sprint6}

\begin{table}[htbp]
\centering
\caption{Backlog du Sprint 6}
\label{tab:backlog_sprint6}
\renewcommand{\arraystretch}{1.3}
\begin{tabular}{|c|p{10cm}|c|}
\hline
\textbf{ID} & \textbf{Tâches} & \textbf{Estimation} \\
\hline
S6-T1 & Définir la stratégie et les scénarios de test du projet & 2 jours \\
\hline
S6-T2 & Exécuter les tests unitaires et d'intégration du backend & 4 jours \\
\hline
S6-T3 & Vérifier les contrôleurs REST, la sécurité et l'authentification & 3 jours \\
\hline
S6-T4 & Tester les processus BPMN, l'assistant RAG et les services métier & 3 jours \\
\hline
S6-T5 & Exécuter les tests frontend et les tests des parcours utilisateur & 3 jours \\
\hline
S6-T6 & Corriger les anomalies et effectuer les régressions nécessaires & 4 jours \\
\hline
S6-T7 & Finaliser la documentation technique et fonctionnelle & 2 jours \\
\hline
S6-T8 & Préparer la livraison et la soutenance du projet & 2 jours \\
\hline
\multicolumn{2}{|r|}{\textbf{Total}} & \textbf{23 jours} \\
\hline
\end{tabular}
\end{table}

\section{Stratégie de test}
\label{sec:strategie_tests_sprint6}

La stratégie de test est organisée en plusieurs niveaux. Les tests
unitaires vérifient le comportement d'une classe ou d'un service de
manière isolée. Les tests d'intégration vérifient la collaboration entre
les composants du backend. Les tests des contrôleurs valident les
réponses HTTP, les règles d'accès et les données retournées par l'API.
Enfin, les tests frontend vérifient le chargement des pages et les
parcours utilisateur les plus importants.

\begin{figure}[htbp]
\centering
\IfFileExists{Figures/test_strategy_sprint6.png}{%
  \includegraphics[width=\textwidth,height=0.75\textheight,keepaspectratio]{Figures/test_strategy_sprint6.png}%
}{%
  \fbox{\parbox[c][5cm][c]{0.88\textwidth}{\centering
  Espace réservé au schéma de la stratégie de test.}}
}
\caption{Stratégie de test de la plateforme FinBank}
\label{fig:test_strategy_sprint6}
\end{figure}

\section{Tests et validation}
\label{sec:tests_validation_sprint6}

\subsection{Tests unitaires du backend}
\label{subsec:tests_unitaires_sprint6}

Les tests unitaires couvrent les services métier et les composants
techniques du backend. Ils concernent notamment l'authentification, la
validation des transactions, la détection d'intention, la recherche
RAG, l'extraction documentaire, les notifications et les traitements
des tâches Flowable.

Chaque test vérifie un comportement précis à partir de données
contrôlées. Les scénarios nominaux sont complétés par des scénarios
d'erreur, par exemple un montant invalide, une décision Flowable
interdite ou l'indisponibilité d'un fournisseur d'intelligence
artificielle.

\begin{figure}[htbp]
\centering
\IfFileExists{Figures/backend_tests_sprint6.png}{%
  \includegraphics[width=0.9\textwidth]{Figures/backend_tests_sprint6.png}%
}{%
  \fbox{\parbox[c][4cm][c]{0.85\textwidth}{\centering
  Espace réservé à la capture de l'exécution des tests backend.}}
}
\caption{Exécution des tests unitaires et d'intégration du backend}
\label{fig:backend_tests_sprint6}
\end{figure}

\subsection{Tests des contrôleurs et de la sécurité}
\label{subsec:tests_controleurs_sprint6}

Les tests des contrôleurs vérifient l'accès aux endpoints et la forme
des réponses. Les fonctionnalités d'inscription, de connexion, de
rafraîchissement du jeton, de profil et de double authentification sont
testées avec des cas valides et invalides.

Les règles de sécurité sont également vérifiées afin de confirmer
qu'un utilisateur ne peut pas accéder aux ressources réservées à
l'administrateur et qu'une requête non authentifiée est correctement
refusée. Les opérations bancaires sont contrôlées pour s'assurer
qu'elles restent liées à l'utilisateur authentifié.

\subsection{Tests des processus BPMN et de l'assistant}
\label{subsec:tests_bpmn_rag_sprint6}

Les scénarios Flowable vérifient le démarrage des processus, la
création des tâches, la validation des décisions et la mise à jour de
l'état final. Les scénarios de rejet, d'escalade et de traitement
incomplet sont également pris en compte lorsque le processus le
prévoit.

Pour l'assistant bancaire, les tests vérifient la détection
d'intention, la conservation des messages, la recherche dans la base
RAG et la gestion d'une réponse lorsque les fournisseurs d'IA sont
indisponibles. L'extraction d'informations depuis un texte ou un
document est également vérifiée.

\begin{figure}[htbp]
\centering
\IfFileExists{Figures/workflow_rag_tests_sprint6.png}{%
  \includegraphics[width=0.9\textwidth]{Figures/workflow_rag_tests_sprint6.png}%
}{%
  \fbox{\parbox[c][4cm][c]{0.85\textwidth}{\centering
  Espace réservé à la capture des tests BPMN et RAG.}}
}
\caption{Validation des workflows BPMN et de l'assistant bancaire}
\label{fig:workflow_rag_tests_sprint6}
\end{figure}

\subsection{Tests frontend et parcours utilisateur}
\label{subsec:tests_frontend_sprint6}

Les tests frontend vérifient que l'application est accessible et que
les pages essentielles se chargent correctement. Les parcours
contrôlés comprennent l'affichage de la page de connexion, l'accès au
tableau de bord, la navigation vers les opérations bancaires, la
conversation avec l'assistant et l'accès aux écrans administrateur.

Les tests vérifient également que les routes protégées respectent le
rôle de l'utilisateur et que l'interface affiche les états d'erreur
lorsqu'une requête échoue.

\begin{figure}[htbp]
\centering
\IfFileExists{Figures/frontend_tests_sprint6.png}{%
  \includegraphics[width=0.9\textwidth]{Figures/frontend_tests_sprint6.png}%
}{%
  \fbox{\parbox[c][4cm][c]{0.85\textwidth}{\centering
  Espace réservé à la capture de l'exécution des tests frontend.}}
}
\caption{Tests des interfaces et des parcours utilisateur}
\label{fig:frontend_tests_sprint6}
\end{figure}

\section{Correction des anomalies}
\label{sec:correction_anomalies_sprint6}

Après l'exécution des tests, les anomalies sont classées selon leur
impact sur la sécurité, les fonctionnalités et l'expérience
utilisateur. Les problèmes bloquants sont corrigés en priorité, puis
les erreurs de validation, d'affichage et de navigation sont traitées.

Une nouvelle exécution des scénarios concernés est effectuée après
chaque correction. Cette démarche permet de vérifier que la correction
résout l'anomalie sans introduire de régression dans une fonctionnalité
déjà validée.

\begin{table}[htbp]
\centering
\caption{Catégories de vérification finale}
\label{tab:verification_finale_sprint6}
\renewcommand{\arraystretch}{1.3}
\begin{tabular}{|p{4.2cm}|p{8.8cm}|}
\hline
\textbf{Catégorie} & \textbf{Éléments vérifiés} \\
\hline
Sécurité & Authentification, autorisations, rôles, 2FA et routes protégées. \\
\hline
Fonctionnel & Comptes, transactions, cartes, investissements, assistant et workflows. \\
\hline
Intégration & Communication frontend/backend, persistence et processus Flowable. \\
\hline
Interface & Navigation, affichage des données, messages d'erreur et responsive design. \\
\hline
Documentation & Endpoints REST, diagrammes, instructions d'installation et utilisation. \\
\hline
\end{tabular}
\end{table}

\section{Finalisation et préparation de la livraison}
\label{sec:finalisation_sprint6}

La finalisation consiste à vérifier la cohérence des sources, des
configurations et de la documentation. Les variables nécessaires à
l'exécution sont identifiées, les étapes de démarrage du backend et du
frontend sont documentées et les fonctionnalités principales sont
regroupées dans une version livrable.

Les diagrammes UML et BPMN sont également révisés afin de rester
alignés avec les composants réellement implémentés. Les captures
d'écran, les résultats de tests et les éléments nécessaires à la
soutenance sont ensuite insérés dans le rapport.

\begin{figure}[htbp]
\centering
\IfFileExists{Figures/delivery_sprint6.png}{%
  \includegraphics[width=0.9\textwidth]{Figures/delivery_sprint6.png}%
}{%
  \fbox{\parbox[c][4cm][c]{0.85\textwidth}{\centering
  Espace réservé au schéma de préparation de la livraison finale.}}
}
\caption{Préparation de la livraison de la plateforme}
\label{fig:delivery_sprint6}
\end{figure}

\section{Conclusion générale}
\label{sec:conclusion_sprint6}

Ce chapitre a présenté le dernier sprint du projet FinBank. Les tests
unitaires, d'intégration, de sécurité et frontend ont permis de
contrôler les principaux composants de la plateforme. Les parcours
bancaires, l'assistant conversationnel, les processus BPMN et les
interfaces ont été vérifiés à différents niveaux.

Les corrections et les régressions ont contribué à stabiliser la
solution avant sa livraison. La finalisation de la documentation, des
diagrammes et des captures prépare la présentation du projet et
constitue la dernière étape de sa réalisation.

\section*{Bilan du projet}
\addcontentsline{toc}{section}{Bilan du projet}

Au terme des six sprints, FinBank propose un backend sécurisé, des
services bancaires opérationnels, un assistant enrichi par une base de
connaissances, des workflows BPMN et des interfaces adaptées aux
utilisateurs et aux administrateurs. L'ensemble forme une plateforme
bancaire modulaire dont les composants peuvent évoluer
indépendamment.

